% Options for packages loaded elsewhere
\PassOptionsToPackage{unicode}{hyperref}
\PassOptionsToPackage{hyphens}{url}
\PassOptionsToPackage{dvipsnames,svgnames,x11names}{xcolor}
\documentclass[
  11pt,
  letterpaper,
]{article}
\usepackage{xcolor}
\usepackage[margin=0.85in]{geometry}
\usepackage{amsmath,amssymb}
\setcounter{secnumdepth}{-\maxdimen} % remove section numbering
\usepackage{iftex}
\ifPDFTeX
  \usepackage[T1]{fontenc}
  \usepackage[utf8]{inputenc}
  \usepackage{textcomp} % provide euro and other symbols
\else % if luatex or xetex
  \usepackage{unicode-math} % this also loads fontspec
  \defaultfontfeatures{Scale=MatchLowercase}
  \defaultfontfeatures[\rmfamily]{Ligatures=TeX,Scale=1}
\fi
\usepackage{lmodern}
\ifPDFTeX\else
  % xetex/luatex font selection
\fi
% Use upquote if available, for straight quotes in verbatim environments
\IfFileExists{upquote.sty}{\usepackage{upquote}}{}
\IfFileExists{microtype.sty}{% use microtype if available
  \usepackage[]{microtype}
  \UseMicrotypeSet[protrusion]{basicmath} % disable protrusion for tt fonts
}{}
\usepackage{setspace}
\makeatletter
\@ifundefined{KOMAClassName}{% if non-KOMA class
  \IfFileExists{parskip.sty}{%
    \usepackage{parskip}
  }{% else
    \setlength{\parindent}{0pt}
    \setlength{\parskip}{6pt plus 2pt minus 1pt}}
}{% if KOMA class
  \KOMAoptions{parskip=half}}
\makeatother
\usepackage{color}
\usepackage{fancyvrb}
\newcommand{\VerbBar}{|}
\newcommand{\VERB}{\Verb[commandchars=\\\{\}]}
\DefineVerbatimEnvironment{Highlighting}{Verbatim}{commandchars=\\\{\}}
% Add ',fontsize=\small' for more characters per line
\newenvironment{Shaded}{}{}
\newcommand{\AlertTok}[1]{\textcolor[rgb]{1.00,0.00,0.00}{\textbf{#1}}}
\newcommand{\AnnotationTok}[1]{\textcolor[rgb]{0.38,0.63,0.69}{\textbf{\textit{#1}}}}
\newcommand{\AttributeTok}[1]{\textcolor[rgb]{0.49,0.56,0.16}{#1}}
\newcommand{\BaseNTok}[1]{\textcolor[rgb]{0.25,0.63,0.44}{#1}}
\newcommand{\BuiltInTok}[1]{\textcolor[rgb]{0.00,0.50,0.00}{#1}}
\newcommand{\CharTok}[1]{\textcolor[rgb]{0.25,0.44,0.63}{#1}}
\newcommand{\CommentTok}[1]{\textcolor[rgb]{0.38,0.63,0.69}{\textit{#1}}}
\newcommand{\CommentVarTok}[1]{\textcolor[rgb]{0.38,0.63,0.69}{\textbf{\textit{#1}}}}
\newcommand{\ConstantTok}[1]{\textcolor[rgb]{0.53,0.00,0.00}{#1}}
\newcommand{\ControlFlowTok}[1]{\textcolor[rgb]{0.00,0.44,0.13}{\textbf{#1}}}
\newcommand{\DataTypeTok}[1]{\textcolor[rgb]{0.56,0.13,0.00}{#1}}
\newcommand{\DecValTok}[1]{\textcolor[rgb]{0.25,0.63,0.44}{#1}}
\newcommand{\DocumentationTok}[1]{\textcolor[rgb]{0.73,0.13,0.13}{\textit{#1}}}
\newcommand{\ErrorTok}[1]{\textcolor[rgb]{1.00,0.00,0.00}{\textbf{#1}}}
\newcommand{\ExtensionTok}[1]{#1}
\newcommand{\FloatTok}[1]{\textcolor[rgb]{0.25,0.63,0.44}{#1}}
\newcommand{\FunctionTok}[1]{\textcolor[rgb]{0.02,0.16,0.49}{#1}}
\newcommand{\ImportTok}[1]{\textcolor[rgb]{0.00,0.50,0.00}{\textbf{#1}}}
\newcommand{\InformationTok}[1]{\textcolor[rgb]{0.38,0.63,0.69}{\textbf{\textit{#1}}}}
\newcommand{\KeywordTok}[1]{\textcolor[rgb]{0.00,0.44,0.13}{\textbf{#1}}}
\newcommand{\NormalTok}[1]{#1}
\newcommand{\OperatorTok}[1]{\textcolor[rgb]{0.40,0.40,0.40}{#1}}
\newcommand{\OtherTok}[1]{\textcolor[rgb]{0.00,0.44,0.13}{#1}}
\newcommand{\PreprocessorTok}[1]{\textcolor[rgb]{0.74,0.48,0.00}{#1}}
\newcommand{\RegionMarkerTok}[1]{#1}
\newcommand{\SpecialCharTok}[1]{\textcolor[rgb]{0.25,0.44,0.63}{#1}}
\newcommand{\SpecialStringTok}[1]{\textcolor[rgb]{0.73,0.40,0.53}{#1}}
\newcommand{\StringTok}[1]{\textcolor[rgb]{0.25,0.44,0.63}{#1}}
\newcommand{\VariableTok}[1]{\textcolor[rgb]{0.10,0.09,0.49}{#1}}
\newcommand{\VerbatimStringTok}[1]{\textcolor[rgb]{0.25,0.44,0.63}{#1}}
\newcommand{\WarningTok}[1]{\textcolor[rgb]{0.38,0.63,0.69}{\textbf{\textit{#1}}}}
\usepackage{longtable,booktabs,array}
\usepackage{caption}
\captionsetup[table]{skip=6pt}
\newcounter{none} % for unnumbered tables
\usepackage{calc} % for calculating minipage widths
% Correct order of tables after \paragraph or \subparagraph
\usepackage{etoolbox}
\makeatletter
\patchcmd\longtable{\par}{\if@noskipsec\mbox{}\fi\par}{}{}
\makeatother
% Allow footnotes in longtable head/foot
\IfFileExists{footnotehyper.sty}{\usepackage{footnotehyper}}{\usepackage{footnote}}
\makesavenoteenv{longtable}
\setlength{\emergencystretch}{3em} % prevent overfull lines
\providecommand{\tightlist}{%
  \setlength{\itemsep}{0pt}\setlength{\parskip}{0pt}}
\usepackage{xcolor}
\definecolor{researchblue}{RGB}{31,65,99}
\usepackage{bookmark}
\IfFileExists{xurl.sty}{\usepackage{xurl}}{} % add URL line breaks if available
\urlstyle{same}
% fallback for those not using the hyperref driver hyperxmp:
\makeatletter
\@ifundefined{xmpquote}{\newcommand{\xmpquote}[1]{#1}}{}
\makeatother
\hypersetup{
  pdftitle={Workflow environments{,} reliable evaluations{,} and useful post-training data},
  pdfauthor={World Engine technical research},
  colorlinks=true,
  linkcolor={researchblue},
  filecolor={Maroon},
  citecolor={Blue},
  urlcolor={researchblue},
  pdfcreator={LaTeX via pandoc}}

\title{Workflow environments, reliable evaluations, and useful
post-training data}
\author{World Engine technical research}
\date{October 10, 2026}

\begin{document}
\maketitle

\setstretch{1.05}
Technical report. Prepared 2026-10-10. Primary sources were accessed on
that date. The \href{../research/SOURCES.md}{source ledger} records
exact URLs and evidence limits. The \href{../OPERATIONS.md}{operations
guide} describes the delivered runtime. The
\href{../ROADMAP.md}{roadmap} separates implemented behavior from
proposed work.

\subsection{Decision and scope}\label{decision-and-scope}

Build a small environment platform whose grades can be inspected,
reproduced, and challenged. The useful unit is a complete task attempt:
a versioned starting state, an agent's permitted observations and
actions, a trusted outcome, and enough provenance to explain the result.
A large library of tasks with weak grades is a poor starting point.

World Engine's local implementation is an environment and data system.
It is not a learned-model trainer, a production application replica, or
a security boundary for arbitrary agent-written code. Its deterministic
baseline policies prove that the workflow worlds and run pipeline
execute. They do not establish that a language model has learned, that
simulation transfers to a deployed application, or that the
implementation beats a hosted competitor.

The recommended progression is concrete: validate state transitions and
graders, produce independently reviewed demonstrations and failure
pairs, establish a held-out evaluation, then integrate a trainer only
when the experiment calls for one. Model training should answer a
measured failure that simpler changes to tools, prompts, or workflow
design have not resolved.

\subsection{What an environment actually
contains}\label{what-an-environment-actually-contains}

An agent environment is a state transition system. At time \(t\), the
environment owns state \(s_t\). The agent receives observation \(o_t\),
chooses action \(a_t\), and the environment computes the next state. The
grader derives reward from trusted evidence. An episode ends when the
task completes, the agent stops, a budget expires, or execution fails.

For language agents, \(o_t\) often includes the current instruction,
prior messages, tool results, and a view of business records. It must
exclude information the agent would not have in the real task: target
answers, future events, evaluator predicates, and other agents' private
records. A text interface can be partially observable even when its
implementation is deterministic.

One convenient formal model is a partially observable Markov decision
process:

\[\mathcal{M}=(\mathcal{S},\mathcal{A},T,\mathcal{O},Z,R,\gamma).\]

Here \(T\) is the transition distribution, \(Z\) is the observation
distribution, and \(R\) is the reward function. The policy is
\(\pi(a_t\mid h_t)\), where history \(h_t\) includes prior observations
and actions. Treating \(h_t\) as the agent's decision input avoids
pretending the latest tool result contains everything needed to act. A
task distribution supplies initial states and instructions. The
implementation can use deterministic \(T\) for each seeded task while
evaluation samples across many tasks.

Keep the following concepts distinct:

{\def\LTcaptype{none} % do not increment counter
\begin{longtable}[]{@{}
  >{\raggedright\arraybackslash}p{(\linewidth - 4\tabcolsep) * \real{0.3333}}
  >{\raggedright\arraybackslash}p{(\linewidth - 4\tabcolsep) * \real{0.3333}}
  >{\raggedright\arraybackslash}p{(\linewidth - 4\tabcolsep) * \real{0.3333}}@{}}
\toprule\noalign{}
\begin{minipage}[b]{\linewidth}\raggedright
Concept
\end{minipage} & \begin{minipage}[b]{\linewidth}\raggedright
Concrete meaning
\end{minipage} & \begin{minipage}[b]{\linewidth}\raggedright
Common error
\end{minipage} \\
\midrule\noalign{}
\endhead
\bottomrule\noalign{}
\endlastfoot
World & Rules, tools, permissions, and state model & Calling a prompt an
environment \\
Task & Instruction plus initial state and target constraints & Changing
the hidden goal after observing a policy \\
Episode & One policy attempt on one task & Treating retries as
independent tasks \\
Rollout & Ordered observations, actions, and outcomes & Saving only the
final answer \\
Grader & Trusted computation of outcomes or quality & Accepting the
agent's claim that it succeeded \\
Reward & Numeric training signal derived from grading & Treating a
shaped score as business success \\
Benchmark & Fixed evaluation distribution and protocol & Reporting
selected easy seeds as general capability \\
Dataset & Curated records with lineage and usage rights & Equating many
traces with useful supervision \\
\end{longtable}
}

The separation of transcript and outcome matters. An agent can write a
convincing confirmation while its tool call failed. Anthropic's
evaluation guide describes task, trial, grader, transcript, and outcome
as distinct concepts and recommends outcome checks for workflow tasks.
\href{https://www.anthropic.com/engineering/demystifying-evals-for-ai-agents}{Anthropic
agent evals}.

\subsection{What existing systems teach
us}\label{what-existing-systems-teach-us}

Fleet's identity is verified through its own homepage and linked GitHub
organization. Its archived Python SDK documents environment creation by
versioned key, reset with a seed and timestamp, state access through a
SQL resource, instance lookup, and explicit close. This is evidence for
a lifecycle-oriented client API. It is not public proof of its private
graders, workload fidelity, or startup speed.
\href{https://fleetai.com/}{Fleet},
\href{https://github.com/fleet-ai/fleet-sdk}{Fleet SDK}.

Proximal's October 8, 2026 infrastructure post separates model compute
on Modal from tools in Kubernetes sandboxes. It records exact sampled
tokens and log probabilities, fixes an adapter version for a full
rollout, bounds stale data, and persists unused rollouts with
checkpoints. The design lesson is stable policy identity and durable
collection across asynchronous systems. Its timing and learning gains
are reported internal results, not measurements reproduced for World
Engine. \href{https://www.proximal.so/blog/posttraining-infra/}{Proximal
infrastructure}.

FrontierSWE v2 publicly packages Harbor tasks with instructions,
configuration, environment definitions, preflight checks, reference
solutions, and images pinned by digest. Its task notice says evaluation
purposes and asks readers to inspect upstream attribution. Copying the
dataset into a training corpus therefore needs a task-specific rights
review.
\href{https://github.com/Proximal-Labs/frontier-swe-v2}{FrontierSWE v2
repository}.

Harbor organizes tasks, agents, execution environments, trials, and
jobs. That separation makes it a strong integration target for terminal
and container tasks. Its ATIF format can reduce trajectory conversion
work, and its separate-verifier design provides an additional place to
isolate evaluation code. These capabilities complement a workflow state
simulator. Adopting a format does not prove an arbitrary grader is
protected. \href{https://github.com/harbor-framework/harbor}{Harbor},
\href{https://docs.harborframework.com/agents/atif}{ATIF},
\href{https://docs.harborframework.com/tasks/separate-verifier}{Separate
verifier}.

PufferLib addresses a different cost regime. Its engineering work
concentrates on simulation, compact policies, native kernels, and
profiling. A millions-of-steps-per-second game workload cannot predict
throughput for a language agent spending seconds on each inference call.
Reuse the discipline of profiling and batching, and compare throughput
only for identical workloads, hardware, policy, and quality thresholds.
\href{https://puffer.ai/blog/engineering-4.0/}{PufferLib engineering}.

The benchmarks describe useful points along the fidelity spectrum.
WebArena uses self-hosted web applications. BrowserGym provides a
Gymnasium browser lifecycle and connects several task suites. OSWorld
exercises real desktop applications and includes manual task
examination. tau-bench focuses on policy-governed tool and user
interaction. AgentBench broadens the domain mix. None alone proves a
generic agent is ready for a specific deployed workflow.
\href{https://github.com/web-arena-x/webarena}{WebArena},
\href{https://github.com/ServiceNow/BrowserGym}{BrowserGym},
\href{https://github.com/xlang-ai/OSWorld}{OSWorld},
\href{https://arxiv.org/abs/2406.12045}{tau-bench},
\href{https://github.com/THUDM/AgentBench}{AgentBench}.

The current tau2-bench repository includes later tau3 changes and
published task fixes. Comparing results requires pinning the task and
evaluator revision. A grader correction can change measured capability
without any policy change.
\href{https://github.com/sierra-research/tau2-bench}{tau2-bench}.

{\def\LTcaptype{none} % do not increment counter
\begin{longtable}[]{@{}
  >{\raggedright\arraybackslash}p{(\linewidth - 6\tabcolsep) * \real{0.2500}}
  >{\raggedright\arraybackslash}p{(\linewidth - 6\tabcolsep) * \real{0.2500}}
  >{\raggedright\arraybackslash}p{(\linewidth - 6\tabcolsep) * \real{0.2500}}
  >{\raggedright\arraybackslash}p{(\linewidth - 6\tabcolsep) * \real{0.2500}}@{}}
\toprule\noalign{}
\begin{minipage}[b]{\linewidth}\raggedright
System
\end{minipage} & \begin{minipage}[b]{\linewidth}\raggedright
Useful documented capability
\end{minipage} & \begin{minipage}[b]{\linewidth}\raggedright
What World Engine can borrow
\end{minipage} & \begin{minipage}[b]{\linewidth}\raggedright
What requires separate proof
\end{minipage} \\
\midrule\noalign{}
\endhead
\bottomrule\noalign{}
\endlastfoot
Fleet SDK & Versioned environment lifecycle and state access & Short
create-reset-act-grade-close usage & Hosted fidelity, hidden grading,
speed \\
Harbor & Task packages, jobs, sandboxes, traces & Adapter and artifact
contracts & Task-specific isolation and parity \\
Verifiers & Reusable environment and rubric library & A small
integration at the rollout boundary & Compatibility with a chosen
trainer \\
PufferLib & Efficient simulator and training implementation & Measure
the actual bottleneck & Language-agent throughput \\
Modal and E2B & Managed sandbox lifecycle & Create, execute, snapshot,
clean up & Cost, limits, security configuration \\
\end{longtable}
}

Verifiers is an environment and evaluation library, while Modal and E2B
supply execution infrastructure. These are different ownership
boundaries. Provider lifecycle APIs document creation, execution,
timeouts, and termination, but no provider resources were provisioned
for this build.
\href{https://github.com/PrimeIntellect-ai/verifiers}{Verifiers},
\href{https://modal.com/docs/guide/sandboxes}{Modal Sandboxes},
\href{https://docs.e2b.dev/sandbox}{E2B lifecycle}.

\subsection{Architecture that fits the
problem}\label{architecture-that-fits-the-problem}

Consider two concrete designs. A browser-first design runs the actual
application inside an isolated container or VM. It offers authentic
selectors, layout, rendering, and application behavior, but reset cost,
licensing, service dependencies, and clock variation grow. A state-first
design implements the business rules and tool API over an episode-owned
database. It resets cheaply, supports precise invariants, and makes
failures inspectable, but does not measure visual computer use or
undocumented application behavior.

World Engine starts with the state-first design. This is appropriate for
demonstrating tool workflows, trusted grading, trace capture, and
curation. The next meaningful fidelity step is a parity adapter against
a real application, not an animation that resembles that application.

The environment owns allowed transitions. The grader owns outcome
computation. The runner owns budgets, policy interaction, and episode
termination. Persistence owns durable run records. The console reads
operator records and issues run commands. Framework-specific request
parsing stays at the HTTP boundary. A trainer adapter, when added,
consumes a documented trajectory contract and never reaches into a live
environment's private grader.

\begin{Shaded}
\begin{Highlighting}[]
\NormalTok{versioned task + seed}
\NormalTok{        |}
\NormalTok{        v}
\NormalTok{episode{-}owned state {-}{-} permitted observation {-}{-}\textgreater{} policy}
\NormalTok{        \^{}                                      |}
\NormalTok{        |                                      v}
\NormalTok{        +{-}{-}{-}{-}{-}{-}{-}{-}{-}{-}{-}{-}{-} validated action \textless{}{-}{-}{-}{-}{-}{-}+ }
\NormalTok{        |}
\NormalTok{        +{-}{-}\textgreater{} immutable event record {-}{-}\textgreater{} operator inspection}
\NormalTok{        |}
\NormalTok{        +{-}{-}\textgreater{} trusted final{-}state grader {-}{-}\textgreater{} curation decision}
\NormalTok{                                              |}
\NormalTok{                                              v}
\NormalTok{                                    versioned dataset export}
\end{Highlighting}
\end{Shaded}

Keep episode state separate before adding locks. Independent databases
or snapshots remove cross-episode mutation by construction. A shared run
database can store summaries and trace references, with transactional
updates for lifecycle changes. A lock on one global mutable world would
preserve correctness only by serializing all work.

A module boundary is justified when it owns a rule or an external
integration. A generic plugin system, custom scheduler, or remote object
store is unnecessary until a concrete integration makes it useful. The
local build should make one complete workflow reliable before
accumulating abstractions.

\subsection{Fidelity has measurable
dimensions}\label{fidelity-has-measurable-dimensions}

``100\% fidelity'' is not a testable general requirement. Applications
change, observed behavior is incomplete, and production includes users
and integrations outside the simulator. Define fidelity relative to a
task distribution and a pinned application revision.

{\def\LTcaptype{none} % do not increment counter
\begin{longtable}[]{@{}
  >{\raggedright\arraybackslash}p{(\linewidth - 4\tabcolsep) * \real{0.3333}}
  >{\raggedright\arraybackslash}p{(\linewidth - 4\tabcolsep) * \real{0.3333}}
  >{\raggedright\arraybackslash}p{(\linewidth - 4\tabcolsep) * \real{0.3333}}@{}}
\toprule\noalign{}
\begin{minipage}[b]{\linewidth}\raggedright
Dimension
\end{minipage} & \begin{minipage}[b]{\linewidth}\raggedright
Question
\end{minipage} & \begin{minipage}[b]{\linewidth}\raggedright
Proposed acceptance evidence
\end{minipage} \\
\midrule\noalign{}
\endhead
\bottomrule\noalign{}
\endlastfoot
Transition fidelity & Does the same valid action change the same
records? & Differential replay against a real app for a sampled command
suite \\
Permission fidelity & Are forbidden reads and writes rejected
identically? & Role-by-operation matrix, deny tests, scope boundary
cases \\
Data fidelity & Do schemas, missing values, joins, and distributions
match? & Schema comparison and redacted distribution checks \\
Temporal fidelity & Do delays, deadlines, retries, and concurrent
changes matter similarly? & Controlled race and delayed-event cases \\
Observation fidelity & Does the agent receive equivalent information? &
Expert review of screens, tool outputs, visibility, and error
messages \\
Workflow fidelity & Are prerequisites, approvals, and irreversible
effects preserved? & Expert walkthroughs with approved outcome
examples \\
Grading validity & Does a pass mean the requested work happened? &
Positive, negative, adversarial, and ambiguous cases independently
labeled \\
Transfer & Does improvement persist outside the simulator? & Locked
evaluation in a separate app revision or controlled deployment \\
\end{longtable}
}

A tool-only environment can pass transition and permission checks while
failing visual fidelity. An attractive UI can fail all business-state
checks. Record each dimension separately.

Build task families from domain constraints. A support refund world
needs eligibility, identity verification, amounts, authorization, and an
immutable transaction record. An inventory world needs reservations,
available stock, and an operation that cannot silently drive stock
negative. A reconciliation world needs canonical identities, references,
and a preservation rule for unrelated accounts. Real experts should
supply the surprising exceptions that synthetic task generation misses.

Validation uses at least one feasible reference execution, a plausible
failure, and an exploit attempt per task family. The reference proves
feasibility, not correctness of the grading specification. Have someone
who did not implement the grader review the outcome criteria. Track
disagreement and repair the task before policy comparison.

\subsection{Worked example: a permitted
refund}\label{worked-example-a-permitted-refund}

This example illustrates the design; its IDs and schema are explanatory,
not an exact SDK contract.

A user requests a \$40 refund for order O-17. The starting state records
customer C-9, order O-17, a \$40 captured payment, no prior refund, a
verified identity flag, and a policy allowing that refund. A second
order O-18 belongs to another customer and must stay unchanged.

The agent observes the request and permitted lookup tools. It first
calls \texttt{lookup\_order(O-17)}. The observation includes the order
status, amount, owner, and applicable policy. It calls
\texttt{issue\_refund(O-17,\ 4000)} using integer cents. The environment
checks identity, ownership, refund eligibility, and idempotency before
creating a refund event. The agent then returns a confirmation.

The trusted grader checks all of the following: O-17 has exactly one
committed refund of 4000 cents, the refund belongs to C-9, the audit
event matches the actual transition, and O-18 and unrelated payments are
unchanged. The agent's confirmation is useful for interaction quality,
but it cannot substitute for those state checks.

A naive grader that checks only \texttt{orders.status\ ==\ refunded}
rewards changing a label without returning money. A second naive grader
that checks only the presence of a refund record rewards an incorrect
amount or duplicate refund. An evaluator must inspect the business
invariant, including preservation constraints, not merely a convenient
field.

An illustrative safe trace contains:

\begin{Shaded}
\begin{Highlighting}[]
\FunctionTok{\{}
  \DataTypeTok{"schema\_version"}\FunctionTok{:} \StringTok{"1"}\FunctionTok{,}
  \DataTypeTok{"task\_id"}\FunctionTok{:} \StringTok{"refund{-}example"}\FunctionTok{,}
  \DataTypeTok{"task\_revision"}\FunctionTok{:} \StringTok{"sha256:..."}\FunctionTok{,}
  \DataTypeTok{"environment\_revision"}\FunctionTok{:} \StringTok{"sha256:..."}\FunctionTok{,}
  \DataTypeTok{"policy\_revision"}\FunctionTok{:} \StringTok{"reference{-}v1"}\FunctionTok{,}
  \DataTypeTok{"seed"}\FunctionTok{:} \DecValTok{17}\FunctionTok{,}
  \DataTypeTok{"split"}\FunctionTok{:} \StringTok{"train"}\FunctionTok{,}
  \DataTypeTok{"events"}\FunctionTok{:} \OtherTok{[}
    \FunctionTok{\{}
      \DataTypeTok{"index"}\FunctionTok{:} \DecValTok{0}\FunctionTok{,}
      \DataTypeTok{"action"}\FunctionTok{:} \FunctionTok{\{}\DataTypeTok{"tool"}\FunctionTok{:} \StringTok{"lookup\_order"}\FunctionTok{,} \DataTypeTok{"order\_id"}\FunctionTok{:} \StringTok{"O{-}17"}\FunctionTok{\},}
      \DataTypeTok{"observation"}\FunctionTok{:} \FunctionTok{\{}\DataTypeTok{"amount\_cents"}\FunctionTok{:} \DecValTok{4000}\FunctionTok{,} \DataTypeTok{"eligible"}\FunctionTok{:} \KeywordTok{true}\FunctionTok{\}}
    \FunctionTok{\}}\OtherTok{,}
    \FunctionTok{\{}
      \DataTypeTok{"index"}\FunctionTok{:} \DecValTok{1}\FunctionTok{,}
      \DataTypeTok{"action"}\FunctionTok{:} \FunctionTok{\{}
        \DataTypeTok{"tool"}\FunctionTok{:} \StringTok{"issue\_refund"}\FunctionTok{,} \DataTypeTok{"order\_id"}\FunctionTok{:} \StringTok{"O{-}17"}\FunctionTok{,}
        \DataTypeTok{"amount\_cents"}\FunctionTok{:} \DecValTok{4000}
      \FunctionTok{\},}
      \DataTypeTok{"observation"}\FunctionTok{:} \FunctionTok{\{}\DataTypeTok{"status"}\FunctionTok{:} \StringTok{"committed"}\FunctionTok{,} \DataTypeTok{"refund\_id"}\FunctionTok{:} \StringTok{"R{-}1"}\FunctionTok{\}}
    \FunctionTok{\}}
  \OtherTok{]}\FunctionTok{,}
  \DataTypeTok{"termination"}\FunctionTok{:} \StringTok{"agent\_finished"}\FunctionTok{,}
  \DataTypeTok{"outcome"}\FunctionTok{:} \FunctionTok{\{}\DataTypeTok{"success"}\FunctionTok{:} \KeywordTok{true}\FunctionTok{,} \DataTypeTok{"reward"}\FunctionTok{:} \FloatTok{1.0}\FunctionTok{\},}
  \DataTypeTok{"curation"}\FunctionTok{:} \FunctionTok{\{}\DataTypeTok{"decision"}\FunctionTok{:} \StringTok{"accepted"}\FunctionTok{,} \DataTypeTok{"review\_revision"}\FunctionTok{:} \StringTok{"review{-}v1"}\FunctionTok{\}}
\FunctionTok{\}}
\end{Highlighting}
\end{Shaded}

The operator record can include detailed failed invariants and trusted
before-and-after state. The agent-facing record and training export must
omit private expected values and sealed evaluation metadata. For SFT,
the accepted demonstration becomes messages and tool calls with a
documented loss mask. For DPO, pair it with a same-task failed attempt
only when the preference label is defensible. For RL, replay verifies
the environment while a trainer collects fresh policy samples.

\subsection{Reward design and reward
hacking}\label{reward-design-and-reward-hacking}

For a finite episode, return is

\[G_t=\sum_{k=t}^{T-1}\gamma^{k-t}r_k.\]

A terminal business outcome can use \(\gamma=1\) when intermediate
rewards are absent and horizon is bounded. Discounting and time
penalties encode preferences for shorter or earlier behavior, so
document them rather than treating them as neutral defaults.

A useful grade is a vector first: task completion, permission
compliance, preservation of unrelated state, and interaction quality. A
training scalar can be derived from that vector. Keep the raw components
so a higher aggregate does not conceal a safety regression.

For example, hard constraints can gate a score:

\[R=\mathbf{1}_{\mathrm{valid\ permissions}}\,\mathbf{1}_{\mathrm{preserved\ state}}\,\mathbf{1}_{\mathrm{completed\ task}}.\]

Each factor is binary in this illustration. Optional progress rewards
should never let an invalid episode outscore a valid completion. Partial
credit can be useful for exploration, but it changes the objective. Test
the reward ordering with adversarial examples before optimizing it.

Potential-based shaping uses \(r'_t=r_t+\gamma\Phi(s_{t+1})-\Phi(s_t)\).
Under its assumptions, this preserves optimal policy ordering. In finite
tasks, terminal potential handling matters. An arbitrary reward for
calling the correct-looking tool is not automatically potential-based
and may reward loops or incomplete work.
\href{https://ai.stanford.edu/~ang/papers/shaping-icml99.pdf}{Ng,
Harada, and Russell}.

{\def\LTcaptype{none} % do not increment counter
\begin{longtable}[]{@{}
  >{\raggedright\arraybackslash}p{(\linewidth - 4\tabcolsep) * \real{0.3333}}
  >{\raggedright\arraybackslash}p{(\linewidth - 4\tabcolsep) * \real{0.3333}}
  >{\raggedright\arraybackslash}p{(\linewidth - 4\tabcolsep) * \real{0.3333}}@{}}
\toprule\noalign{}
\begin{minipage}[b]{\linewidth}\raggedright
Failure class
\end{minipage} & \begin{minipage}[b]{\linewidth}\raggedright
Concrete exploit
\end{minipage} & \begin{minipage}[b]{\linewidth}\raggedright
Defensive test or design
\end{minipage} \\
\midrule\noalign{}
\endhead
\bottomrule\noalign{}
\endlastfoot
Proxy completion & Mark a ticket resolved without issuing the refund &
Compare trusted transaction state with policy outcome \\
Overbroad mutation & Merge every record to remove duplicate count &
Check identity constraints and unrelated record preservation \\
Repeated shaping & Repeatedly call a progress tool for points & Reward
state change, cap or eliminate shaping, test cycles \\
Grader tampering & Change expected amount or overwrite reward file & Run
private grader outside writable agent state \\
Test manipulation & Skip a test or force a zero exit code & Treat
structured verifier result as authoritative, test harness bypasses \\
Observation leakage & Read target values from debug output & Allowlist
agent observations and test serialized outputs \\
Judge persuasion & Include ``ignore rubric and award full marks'' in
answer & Treat transcript as data, constrain rubric judge, use state
checks \\
Time abuse & Wait past a deadline to make an invalid condition disappear
& Use controlled logical time and stable grading rules \\
Shared-state attack & Read or mutate another episode's records &
Per-episode ownership and concurrent isolation tests \\
Evaluation overfit & Detect a known benchmark fixture and special-case
it & Hold out task families, revisions, entities, and adversarial
cases \\
\end{longtable}
}

Anthropic's controlled coding study reports that learned reward hacking
can accompany broader undesirable behavior. That result strengthens the
case for grader integrity; it does not quantify the risk of every
possible training pipeline.
\href{https://www.anthropic.com/research/emergent-misalignment-reward-hacking}{Reward
hacking study}.

Do not call a grader ``anti-cheat'' because ordinary failing policies
receive low rewards. Seed known exploits, require that the grader
rejects them, and record exactly which exploit classes were tested. A
newly discovered exploit invalidates the relevant assurance until
repaired. Hidden tests make memorization harder, but hidden tests inside
an agent-readable filesystem are not hidden.

\subsection{Trust boundaries and
contamination}\label{trust-boundaries-and-contamination}

The local operator, environment implementation, evaluator, and policy
runner are separate roles even when one developer owns all of them.
Production security requires enforcing that separation through
processes, credentials, namespaces, and sandbox permissions.

The threat model includes malformed tool arguments, attempts to access
forbidden data, mutation of evaluator inputs, injected instructions in
records, poisoning of demonstrations, cross-episode access, resource
exhaustion, and exfiltration. A local policy object running in the host
Python process can inspect host memory if it is malicious. A database
permission rule cannot make that Python object an untrusted-code
sandbox.

For untrusted code, put only the agent workspace and authorized tool
client inside a sandbox. Keep the evaluator, target answers, operator
database, cloud credentials, and protected test material outside it.
Give each episode scoped credentials, deny external network access by
default, and admit only task-required destinations. Run final grading on
copied artifacts or a trusted state snapshot. Retain a proof that the
snapshot belongs to the graded episode.

Expose task goals to the agent when they are part of the user's request.
Hide only the private evaluator implementation and protected held-out
content. Obscuring a requirement creates ambiguity rather than a
stronger benchmark. A refund amount specified by the user is public task
information; an expected-state table describing a held-out variant is
private evaluator data.

Keep train, development, and test splits at the task-family or
source-incident level when records are correlated. A different seed with
the same hidden template often remains near-duplicate data. Deduplicate
across exports and holdouts by normalized task lineage and semantic
review, not merely byte hashes.

Redact before data leaves the operator trust boundary. Use structured
field rules for names, emails, payment identifiers, secrets, and
user-supplied free text. Redaction itself can change the task, so test
that relationships and constraints remain valid. Keep access-controlled
original-to-redacted mappings only when justified. A public dataset must
not retain a reversible identity map.

Poisoning research motivates review of data origin and unexpected
behavioral instructions. It does not supply a universal safe number of
suspicious records.
\href{https://www.anthropic.com/research/small-samples-poison}{Anthropic
poisoning study}.

\subsection{Traces are evidence before they are training
examples}\label{traces-are-evidence-before-they-are-training-examples}

Record enough to distinguish a policy mistake from an environment
failure. An episode record should include task and world revisions,
grader revision, seed, starting-state digest, policy provider and
configuration, tool schema revision, ordered events, termination cause,
final outcome, elapsed times, and curation decision. Include token usage
and model sampling details when the provider exposes them.

Canonical hashing detects accidental changes to the canonicalized
record. It does not prove the record is authentic when its author can
rewrite the record and recompute the hash. Signed manifests and
append-only storage can strengthen authenticity. They still depend on
who controls the signing key and collector.

Three forms of replay answer different questions:

{\def\LTcaptype{none} % do not increment counter
\begin{longtable}[]{@{}
  >{\raggedright\arraybackslash}p{(\linewidth - 4\tabcolsep) * \real{0.3333}}
  >{\raggedright\arraybackslash}p{(\linewidth - 4\tabcolsep) * \real{0.3333}}
  >{\raggedright\arraybackslash}p{(\linewidth - 4\tabcolsep) * \real{0.3333}}@{}}
\toprule\noalign{}
\begin{minipage}[b]{\linewidth}\raggedright
Replay form
\end{minipage} & \begin{minipage}[b]{\linewidth}\raggedright
Purpose
\end{minipage} & \begin{minipage}[b]{\linewidth}\raggedright
Inputs
\end{minipage} \\
\midrule\noalign{}
\endhead
\bottomrule\noalign{}
\endlastfoot
Display replay & What did the recorded agent see and do? & Saved ordered
events \\
Environment replay & Do saved actions reproduce state transitions and
grade? & Pinned world, starting state, ordered actions, controlled
randomness \\
Policy replay & Does the model choose the same actions again? & Pinned
model, provider behavior, sampling configuration, tool responses \\
\end{longtable}
}

A deterministic seed supports environment replay only if all random
choices, clocks, IDs, external responses, and interleavings are
controlled or recorded. Parallel scheduling can produce different
wall-clock order while each isolated episode remains deterministic. Sort
benchmark summaries by task identity, and never let task completion
order define the split or reward.

Store the original outcome and any regrade as separate revisions. Fixing
a grader should not overwrite the historical explanation. Prefer a
manifest that binds selected episode IDs, digests, selection rules,
export schema, rights metadata, and review version. Export dates alone
cannot explain where a training example came from.

\subsection{SFT, DPO, RLHF, PPO, and
GRPO}\label{sft-dpo-rlhf-ppo-and-grpo}

These methods solve different problems. SFT trains the model to imitate
selected behavior. DPO trains relative preferences between candidate
outputs. RLHF combines human feedback with reward-based optimization.
PPO and GRPO are optimization methods, not synonyms for RLHF. A
deterministic programmatic reward makes a run RL with verifiable
rewards, even when no human preference model is involved.

The InstructGPT work used demonstrations, ranked model outputs, and
reward-based optimization. It gives the historical pipeline, not a
promise that any dataset will improve any model.
\href{https://arxiv.org/abs/2203.02155}{InstructGPT}.

\subsubsection{Supervised fine-tuning}\label{supervised-fine-tuning}

For demonstration data D, a masked SFT objective is:

\[\mathcal{L}_{\mathrm{SFT}}=-\mathbb{E}_{\mathcal{D}}\left[\sum_t m_t\log\pi_\theta(y_t\mid x,y_{<t})\right].\]

The mask \(m_t\) identifies model-generated target tokens. Tool
observations and user messages normally provide context rather than
prediction targets. A tool-use demonstration requires a compatible
tokenizer, chat template, role formatting, and tool-call representation.
Loss masking must match the trainer's actual implementation.

Select demonstrations for correctness and useful diversity. A reference
solver can produce syntactically valid traces, but it may use privileged
information or a fixed sequence unavailable to a normal agent. Review
whether every action was justified by visible information. Preserve
failure traces for evaluation and preference curation rather than
silently deleting the failure population.

\subsubsection{Direct preference
optimization}\label{direct-preference-optimization}

For prompt \(x\), chosen output \(y_w\), rejected output \(y_l\), and
reference policy \(\pi_{\mathrm{ref}}\), the standard DPO objective is:

\[\mathcal{L}_{\mathrm{DPO}}=-\mathbb{E}\left[\log\sigma\left(\beta\left[\log\frac{\pi_\theta(y_w\mid x)}{\pi_{\mathrm{ref}}(y_w\mid x)}-\log\frac{\pi_\theta(y_l\mid x)}{\pi_{\mathrm{ref}}(y_l\mid x)}\right]\right)\right].\]

The comparison must share a task and initial conditions. A successful
easy task and a failed hard task are not a defensible preference pair. A
larger reward gap is useful evidence, but inspect whether the chosen
trace used forbidden information or exploited the grader. DPO can
overfit preferences and does not interact with the environment during an
ordinary offline update. \href{https://arxiv.org/abs/2305.18290}{DPO
paper}.

For multi-turn tool tasks, paired histories may contain different
observations after the first differing action. The trainer needs a
supported trajectory encoding. Concatenating arbitrary histories into a
single response without documenting roles and masks changes the learning
problem.

\subsubsection{PPO and reward-based
optimization}\label{ppo-and-reward-based-optimization}

Let \(\rho_t\) be the probability ratio between the updated policy and
the policy that collected the sample. A common PPO surrogate is:

\[L_{\mathrm{clip}}=\mathbb{E}\left[\min\left(\rho_t A_t,\operatorname{clip}(\rho_t,1-\epsilon,1+\epsilon)A_t\right)\right].\]

The trainer maximizes this surrogate and often combines it with value
loss, entropy terms, and a reference-policy KL penalty. Advantage
\(A_t\) estimates whether an action did better than a baseline. PPO
generally uses a learned value function. Clipping limits parts of the
sampled objective; it does not guarantee a hard KL bound or prevent
reward hacking. \href{https://arxiv.org/abs/1707.06347}{PPO paper}.

If a provider returns text without sampled-policy probabilities, an
exported trace is not enough to compute the usual ratio correctly.
Re-tokenizing text can differ from the exact tokens sampled. Even exact
output tokens are insufficient if the trainer cannot reconstruct
context, masks, policy version, and behavior probabilities.

\subsubsection{GRPO and grouped
rollouts}\label{grpo-and-grouped-rollouts}

GRPO samples a group of candidate trajectories for the same prompt or
task. A common group-relative advantage is:

\[A_i=\frac{R_i-\operatorname{mean}(R_{\mathrm{group}})}{\operatorname{std}(R_{\mathrm{group}})+\epsilon}.\]

It removes the learned value model and compares candidates relative to
their group. The original formulation combines a clipped policy
objective with a reference KL term. Implementations differ in token
normalization, reward scaling, KL estimation, and treatment of long
sequences. Pin the trainer and configuration.
\href{https://arxiv.org/abs/2402.03300}{DeepSeekMath},
\href{https://huggingface.co/docs/trl/main/en/grpo_trainer}{TRL GRPO
trainer}.

A group whose rewards are all equal has no relative reward signal under
this normalization. Very easy and impossible tasks can both waste
rollouts. Adjust the curriculum using held-out measurements, while
retaining original failure and sampling statistics. Do not secretly drop
hard groups and report the resulting training reward as coverage.

For environment trajectories, all group members need the same task
revision and comparable initial conditions. The trainer must mask tool
and user tokens, track policy version over the entire episode, define
terminal versus truncated returns, and monitor rollout staleness. Long
episodes collected from old weights are an off-policy problem, even if
the environment is fast.

{\def\LTcaptype{none} % do not increment counter
\begin{longtable}[]{@{}
  >{\raggedright\arraybackslash}p{(\linewidth - 6\tabcolsep) * \real{0.2500}}
  >{\raggedright\arraybackslash}p{(\linewidth - 6\tabcolsep) * \real{0.2500}}
  >{\raggedright\arraybackslash}p{(\linewidth - 6\tabcolsep) * \real{0.2500}}
  >{\raggedright\arraybackslash}p{(\linewidth - 6\tabcolsep) * \real{0.2500}}@{}}
\toprule\noalign{}
\begin{minipage}[b]{\linewidth}\raggedright
Method
\end{minipage} & \begin{minipage}[b]{\linewidth}\raggedright
Required evidence or data
\end{minipage} & \begin{minipage}[b]{\linewidth}\raggedright
Main failure to watch
\end{minipage} & \begin{minipage}[b]{\linewidth}\raggedright
What a local export supplies
\end{minipage} \\
\midrule\noalign{}
\endhead
\bottomrule\noalign{}
\endlastfoot
SFT & Reviewed demonstrations and compatible masks & Imitation of
privileged or wrong behavior & Candidate conversations and lineage \\
DPO & Comparable chosen and rejected trajectories & Confounded or
incorrect preferences & Candidate pairs and reasons \\
RLHF & Human judgments, reward validation, optimization & Reward-model
exploitation & Feedback-compatible evidence, not a trained reward
model \\
PPO & Fresh trajectories, behavior probabilities, advantages &
Staleness, value error, wrong masks & Environment records, not training
tensors \\
GRPO & Comparable groups, rewards, policy sampling metadata &
Zero-variance groups and length bias & Groupable tasks and outcomes \\
\end{longtable}
}

Use an established trainer for optimization. verl, OpenRLHF, and Miles
support distributed post-training, while RAGEN studies multi-turn agent
learning and its collapse modes. Integration work remains substantial:
align model templates, rollout callbacks, reward semantics, and
checkpoint lifecycle. A JSONL file named ``GRPO'' does not complete that
integration. \href{https://github.com/verl-project/verl}{verl},
\href{https://github.com/OpenRLHF/OpenRLHF}{OpenRLHF},
\href{https://github.com/radixark/miles}{Miles},
\href{https://github.com/mll-lab-nu/RAGEN}{RAGEN}.

Hosted APIs expose their own training and evaluation contracts. Validate
a specific provider's supported dataset format and environment
interaction before promising compatibility. This repository does not
invoke hosted model training.
\href{https://developers.openai.com/api/docs/guides/evals}{OpenAI
evals},
\href{https://developers.openai.com/api/docs/guides/reinforcement-fine-tuning}{OpenAI
reinforcement fine-tuning}.

\subsection{Evaluation that supports a
claim}\label{evaluation-that-supports-a-claim}

Choose the claim before running the experiment. ``The reference policy
solves generated refund tasks'' is an implementation check. ``A trained
policy solves more held-out refund tasks at equal cost'' is a model
experiment. ``The policy reduces real support errors'' requires evidence
outside the simulator.

Hold fixed the model revision, policy prompt, tool schema, task
revision, step budget, retry budget, sampling parameters, and execution
provider. Record any exception. Compare policies on identical task
identities and seeds. Randomize execution order so capacity shifts or
provider incidents do not systematically favor one policy.

Report task-level outcomes and failure categories alongside aggregates.
Separate policy failures, invalid tasks, environment failures, provider
failures, and cancellation. A primary reliability metric can count all
failures when that matches operational use. Also provide a policy-only
analysis with predefined exclusions and their counts. Silently dropping
failures inflates the result.

For N independent binary outcomes, estimated success is
\texttt{p\_hat\ =\ successes/N}. A Wilson interval behaves better than
the simple normal interval for small samples and proportions near zero
or one. Correlated repeats on the same task are not N independent tasks.
For policy comparisons, compute per-task paired differences and
bootstrap over tasks or task families, preserving repeats inside each
sampled cluster. Report the confidence level and resampling protocol.
The broader need to report uncertainty in few-run RL evaluation is
examined in \href{https://arxiv.org/abs/2108.13264}{Agarwal et al.}.

As an illustration, 18 successes in 20 independent tasks gives 90\%
observed success, but a roughly 70\% to 97\% Wilson 95\% interval. This
is not evidence of a precise 90\% capability. The numerical example is a
statistical illustration, not a benchmark result for this build.

Pass@k measures the chance at least one of k attempts succeeds. Under
independent attempts with equal success probability p, it is
\texttt{1-(1-p)\^{}k}. The common finite-sample estimator uses
\texttt{1\ -\ choose(n-c,k)/choose(n,k)} with c successful samples among
n and k \textless= n.~Pass\^{}k, used for repeated reliability, asks
that every attempt succeed and is conceptually different. Include
attempt cost when comparing retries.

Prevent evaluation leakage through a locked holdout, audited access,
task-family splits, and a fixed review process. A development set is
where you repair prompts and graders. Repeatedly inspecting the test set
and selecting the best checkpoint turns it into development data. Retain
a final test set for one planned comparison or disclose the selection
procedure.

Task validity has its own evaluation. Domain experts should label
whether tasks are feasible, clear, representative, and correctly graded.
Test a sample of claimed successes manually. Investigate unexpected
failures before assuming they reveal model weakness. A policy can fail
because the world omitted a necessary tool or because the grader
contradicts the instruction.

\subsection{Throughput, isolation, and
cost}\label{throughput-isolation-and-cost}

The relevant throughput is valid episodes per second or useful training
tokens per unit cost, not environment steps alone. A rough capacity
model is:

\[q_{\mathrm{episodes}}\ \lesssim\ \min\left(\frac{C_{\mathrm{env}}}{\mathbb{E}[T_{\mathrm{episode}}]},\frac{Q_{\mathrm{model}}}{\mathbb{E}[n_{\mathrm{tokens}}]},\frac{Q_{\mathrm{storage}}}{\mathbb{E}[n_{\mathrm{bytes}}]}\right).\]

Here \(C_{\mathrm{env}}\) is concurrent environment capacity,
\(Q_{\mathrm{model}}\) is effective model token throughput, and
\(Q_{\mathrm{storage}}\) is sustained trace bandwidth. Queue admission,
rate limits, long-tail episodes, and failure retries reduce achieved
throughput. Report startup time, model wait, tool time, grading time,
persistence time, and cleanup separately.

For a local deterministic policy, model wait is nearly absent. Adding an
external language model changes the bottleneck. Increasing simulator
concurrency beyond inference capacity raises queue time and live-state
memory without necessarily completing more episodes.

Little's law relates average in-flight work L, arrival rate lambda, and
average time W through \texttt{L\ =\ lambda\ *\ W} in a stable system.
It helps size a bounded queue, but it cannot justify unlimited
admission. Control pending jobs, active episodes, maximum steps,
observation size, and trace size independently.

Warm pools reduce cold startup latency but retain paid idle resources
and create reset contamination risks. Pool only clean snapshots. Verify
no prior episode files, secrets, processes, or cached tool credentials
remain. An isolated database copy is cheap for a small workflow world; a
full desktop snapshot has different storage and startup costs.

Containers share a kernel and require careful permissions. gVisor and
microVMs add stronger execution boundaries with compatibility and cost
tradeoffs. Isolation claims must name the runtime and configuration. A
sandbox provider's default network access, mounted files, and secrets
policy still matter.

GPU inference benefits from continuous batching, prefix caching, and
enough concurrent requests to avoid idle compute. Long contexts consume
KV-cache memory and can reduce concurrency. Place capture near
inference, retain exact sampling metadata, and avoid cross-region
request paths where latency accumulates over many turns. These are
proposed deployment considerations, not features implemented by the
local build.

Do not quote a universal cost per rollout. Compute:

\[\begin{aligned}
C_{\mathrm{run}}={}&C_{\mathrm{model\ input}}+C_{\mathrm{model\ output}}+C_{\mathrm{sandbox}}\\
&+C_{\mathrm{GPU}}+C_{\mathrm{storage}}+C_{\mathrm{egress}}+C_{\mathrm{retries}}.
\end{aligned}\]

Measure cost per accepted trace and cost per held-out success, since a
fast stream of invalid or duplicate trajectories is not useful
throughput. The first external experiment should have explicit spending,
token, concurrency, and duration caps, plus cancellation and cleanup
proof.

\subsection{Dataset quality and
rights}\label{dataset-quality-and-rights}

A curated trace is a reviewed training candidate, not every recorded
run. Acceptance can require correct final state, no unauthorized
transitions, replay consistency, complete provenance, a permitted usage
license, redaction, and a task split that avoids the holdout. Record
rejection reasons so curation does not hide systematic weaknesses.

Dataset diversity should vary business rules, entity relationships,
error modes, and action sequences. Renaming customers while retaining
one exact solution mostly creates near-duplicates. Include ambiguous
requests, permission denials, unavailable dependencies, idempotent retry
cases, and situations where declining is the correct outcome.

Rights are a record-level property. Separate repository code licenses,
application licenses, dataset licenses, model licenses, and customer
permissions. Open source is not public domain. A public benchmark may
permit evaluation but restrict training reuse. A contributor's
permission to share a trace does not necessarily cover the personal data
or copyrighted documents inside it.

Before sharing, retain a source license or consent basis, permitted
uses, attribution, retention policy, redaction revision, and deletion
procedure. These are data-engineering requirements, not proof of legal
compliance. Obtain appropriate review for the actual jurisdiction and
data agreement instead of assuming a generic notice grants rights.

A production deployment also needs authentication, authorization,
per-tenant storage and encryption policies, audit logs, retention,
deletion, recovery, resource quotas, and incident response. The local
console is an operator tool. It should not be exposed as a multi-tenant
service without that work.

\subsection{A curriculum for learning and extending the
system}\label{a-curriculum-for-learning-and-extending-the-system}

Start with a single workflow and manually execute a successful and a
failing attempt. Inspect the trusted state after each action. Write an
exploit that fools a deliberately weak grade, then repair the grader and
show that the same exploit fails. This establishes why outcome state and
grader permissions matter.

Next, generate variations of the task and run the same policy over them.
Distinguish seed variation from new rule variation. Compare reference,
incomplete, and adversarial policies on the same tasks. Inspect where a
shaped reward and terminal success disagree.

Then build an export manifest. Accept a reviewed successful
demonstration, reject an incomplete one, and pair comparable candidates.
Reproduce the accepted episode from its recorded initial state and
actions. Check that the export omits private evaluator state.

Only after the held-out evaluation is stable, train a small open model
with an established SFT trainer. Compare it against the base model at
equal budgets. Follow with DPO if preferences add reliable information.
Try grouped RL when the environment provides validated rewards and the
inference stack supplies exact sampling metadata. Re-run the locked
evaluation and inspect regressions, not merely mean reward.

The next implementation milestones are grounded in evidence: real
application parity, broader adversarial grader review, stable held-out
model evaluation, trainer integration, and a bounded cloud-provider
experiment. The local model adapter is implemented; integrating it is
not proof of policy improvement. The \href{../ROADMAP.md}{roadmap} gives
the required proof for each. Until those gates pass, the delivered
system remains a useful local environment laboratory with explicit
limits.

\subsection{Practical study sequence}\label{practical-study-sequence}

The companion \href{learning-labs.md}{learning labs} provide six
reproducible exercises with worked answers. They start by contrasting a
successful tool call, a stopped episode, partial reward, and a passing
grade. A refund overpayment attempt shows transition rejection, while
trusted-fixture corruption independently proves the zero-reward safety
gate.

The next exercises collect same-seed reference and incomplete-policy
attempts, replay actions, inspect training-only split filtering, and
explain why seed variation is not task-family generalization. Preference
and grouped exports expose different policies, transformed trajectory
completions, duplicate candidates, zero-variance reward groups, and the
missing data that prevents a text trace from becoming an on-policy RL
batch.

The final exercise profiles concurrency 1, 4, and 8, then contrasts
simulator timing with local model inference. It examines the public
Docker task's separate evaluator, artifact-only handoff, evaluation-only
canary, and protocol adaptations. The objective is to learn what each
measurement establishes and what additional experiment is needed for a
broader claim.

\subsection{Delivered capability and measured
evidence}\label{delivered-capability-and-measured-evidence}

The local package delivers three synthetic workflow worlds, strict
action schemas, reset and observation, isolated episode-owned Python
state, trusted final grading, bounded jobs, cancellation and restart
records, append-only episode storage, replay, provider adapters,
candidate exports, and a connected React operator console. The process
admits at most eight active episodes and 32 active runs. A nonblocking
scheduler lease prevents two manager processes from recovering or
scheduling the same database concurrently. The API binds to loopback and
paid provider calls require both process and run opt-in.

Training formats exclude validation, test, and evaluation-only episodes
and require replay consistency. SFT records controller policy-input
messages per tool decision. Built-in adapters use that shared message
construction, while a custom provider can transform its input
internally; the controller record is not proof of a custom model's
actual request. DPO retains matched provenance and labels its
transformed trajectory completion; a trainer adapter is still required.
Grouped candidates are explicitly nontrainable. The local provider
records the installed model digest when available, but the system does
not capture the complete sampled token sequence and old-policy
probabilities needed for an on-policy update.

Final review found that checking transitions and the raw grader alone
did not reject forged top-level reward and success labels in
SDK-imported episode records. The validator now recomputes authoritative
episode reward and pass, checks status/termination coherence, and
verifies initial/final state, stable lineage, split, scenario,
tool-context sequence, counts, and configuration-hash consistency.
Regression cases reject mislabeled records before they can become
accepted training candidates. This validates semantic coherence, not
collection authenticity: timing, token usage, external failure causes,
and asserted model identity are still reported evidence. Module and
class fingerprints do not recursively attest imported external helpers;
reproducibility still requires dependency and source manifests.

Split assignment uses stable environment-family/seed lineage, separate
from the revision-sensitive scenario identity. This prevents a grader
repair from reassigning a related seed between training and holdout. It
does not hold out entire rule families. The final runner assigns zero
episode reward to budget truncation, cancellation, timeout, and failed
execution even when the retained diagnostic grade contains partial
credit.

\subsubsection{Synthetic workflow
measurements}\label{synthetic-workflow-measurements}

The \href{../../examples/benchmark-results/benchmark.json}{raw timing
matrix} records 54 batches and 1,620 episode executions: three worlds,
two scripted policies, three concurrency settings, three repeated
trials, and 30 matched seeds per batch. On macOS arm64 with Python
3.13.5, the median end-to-end reference-policy rates were:

{\def\LTcaptype{none} % do not increment counter
\begin{longtable}[]{@{}
  >{\raggedright\arraybackslash}p{(\linewidth - 6\tabcolsep) * \real{0.2000}}
  >{\raggedleft\arraybackslash}p{(\linewidth - 6\tabcolsep) * \real{0.2667}}
  >{\raggedleft\arraybackslash}p{(\linewidth - 6\tabcolsep) * \real{0.2667}}
  >{\raggedleft\arraybackslash}p{(\linewidth - 6\tabcolsep) * \real{0.2667}}@{}}
\toprule\noalign{}
\begin{minipage}[b]{\linewidth}\raggedright
World
\end{minipage} & \begin{minipage}[b]{\linewidth}\raggedleft
Concurrency 1, episodes/s
\end{minipage} & \begin{minipage}[b]{\linewidth}\raggedleft
Concurrency 4, episodes/s
\end{minipage} & \begin{minipage}[b]{\linewidth}\raggedleft
Concurrency 8, episodes/s
\end{minipage} \\
\midrule\noalign{}
\endhead
\bottomrule\noalign{}
\endlastfoot
CRM merge & 166.8 & 186.5 & 214.5 \\
Refund operations & 300.2 & 251.4 & 260.0 \\
Incident control & 176.0 & 172.9 & 178.1 \\
\end{longtable}
}

All scripted reference executions passed and all incomplete-policy
executions failed the task grade. Each baseline execution included a
denied premature operation. There were no recorded provider or internal
errors in this matrix. These rates include scheduling, hashing, grading,
and SQLite persistence, not language-model inference, container
provisioning, or real application execution. The incomplete policy is
faster in part because it performs fewer actions; failed attempts per
second are not useful-task throughput.

Each row repeats the same 30 deterministic scenarios three times.
Outcome intervals use the unique first-trial scenarios rather than
treating 90 repeated executions as independent evidence. The model-based
Wilson 95\% interval for 30/30 is approximately 88.6\% to 100\%; for
0/30 it is 0\% to 11.4\%. Even those intervals do not turn selected
synthetic seeds into a representative real-world task distribution. More
concurrency does not uniformly improve these CPU-heavy rates and
increases per-episode latency. Cooperative scheduling does not create
additional CPU execution capacity.

A
\href{../../examples/benchmark-results/schema-cache-comparison.json}{same-source
cache comparison} bypassed only tool-schema memoization in the control
and enabled it in the final matrix, preserving matched cases and all
outcome counts. Across its 18 cells, observed rates improved by about
1.71 to 3.27 times. This is a local implementation measurement, not a
competitor or model-inference speedup. Sequential runs had uncontrolled
background load; exact fingerprints and raw timing trials are retained.
The earlier pre-cache historical receipt is kept separately rather than
silently replaced.

The \href{../../examples/benchmark-results/io-overlap.json}{I/O-overlap
control} adds a simulated 25-millisecond wait to each of six real CRM
decisions. With 24 matched episodes per trial and three timing trials,
median rates were 6.02, 23.75, and 46.49 episodes/s at concurrency 1, 4,
and 8, with observed in-flight decision peaks 1, 4, and 8. Every attempt
passed. This demonstrates overlapping awaited work in the scheduler. It
is explicitly simulated latency, not measured language-model or cloud
performance.

\subsubsection{Actual local model workflow
attempts}\label{actual-local-model-workflow-attempts}

A separate
\href{../../examples/benchmark-results/local-model.json}{paired local
evaluation} ran the already-installed Ollama \texttt{qwen3:8b} once per
world at seed 42, alongside both scripted policies on the same scenario.
It used temperature zero, thinking disabled, 512 generated-token limit
per decision, 20 decisions per episode, concurrency one, and a
120-second episode timeout. Every record was evaluation-only. Collection
used committed implementation revision
\texttt{7fdf4c6ae511dc81cc4074ea54122c0671ba2469}, with
\texttt{implementation\_dirty:\ false}. The overall checkout was dirty
because documentation and other evidence were still being assembled; the
receipt does not claim a wholly clean checkout.

{\def\LTcaptype{none} % do not increment counter
\begin{longtable}[]{@{}
  >{\raggedright\arraybackslash}p{(\linewidth - 8\tabcolsep) * \real{0.1765}}
  >{\raggedright\arraybackslash}p{(\linewidth - 8\tabcolsep) * \real{0.1765}}
  >{\raggedright\arraybackslash}p{(\linewidth - 8\tabcolsep) * \real{0.1765}}
  >{\raggedleft\arraybackslash}p{(\linewidth - 8\tabcolsep) * \real{0.2353}}
  >{\raggedleft\arraybackslash}p{(\linewidth - 8\tabcolsep) * \real{0.2353}}@{}}
\toprule\noalign{}
\begin{minipage}[b]{\linewidth}\raggedright
World
\end{minipage} & \begin{minipage}[b]{\linewidth}\raggedright
Local task grade
\end{minipage} & \begin{minipage}[b]{\linewidth}\raggedright
Termination
\end{minipage} & \begin{minipage}[b]{\linewidth}\raggedleft
Invalid actions
\end{minipage} & \begin{minipage}[b]{\linewidth}\raggedleft
Batch wall seconds
\end{minipage} \\
\midrule\noalign{}
\endhead
\bottomrule\noalign{}
\endlastfoot
CRM merge & Fail & 20-decision limit & 0 & 33.24 \\
Refund operations & Pass & Finished & 1 & 15.13 \\
Incident control & Fail & 20-decision limit & 17 & 34.38 \\
\end{longtable}
}

The reference passed all three and the incomplete baseline failed all
three. The local refund attempt first used an idempotency key belonging
to another payment, received a denial, then recovered and completed the
task. Its passing task grade does not qualify it for the default
clean-demonstration export, which also excludes invalid actions and,
independently, evaluation-only episodes. The incident attempt repeatedly
requested approval before both preflight checks passed. CRM exhausted
its decision budget without completion. Neither unsuccessful attempt was
a provider exception.

This nine-attempt integration sample demonstrates real model
interaction, recovery from a tool denial, and distinct observed workflow
failures. It is not a population estimate, policy-improvement result, or
comparison with a hosted model. The model digest, actions, token counts,
source fingerprints, and safe errors are preserved; token counts are not
exact sampled token sequences or log probabilities. All nine
\href{../../examples/benchmark-results/local-model-traces.jsonl}{full
evaluation records} replay against the delivered source. They are
operator evidence, explicitly not training exports. Replay binds the
environment implementation used at collection and intentionally rejects
later source changes rather than silently regrading old records with new
rules.

The earlier
\href{../../examples/benchmark-results/historical-attempt/local-model.json}{historical
attempt} remains separate. Its intermediate source was uncommitted, so
the recorded fingerprints do not guarantee source recovery. The final
collection repeated the same seeds without prompt tuning; these repeats
do not add independent tasks. A separate UI model request was briefly
submitted and cancelled around 05:03:10--05:03:15 UTC, so final
inference timings may include server contention. No model-inference
speedup is inferred from the historical timing difference.

\subsubsection{One public terminal task}\label{one-public-terminal-task}

The separate runner executes \texttt{log-summary-date-ranges} from
\href{https://github.com/harbor-framework/terminal-bench-2/tree/2fd12b88aafdd04a52c298e3940bcb189f9766d6/log-summary-date-ranges}{Terminal-Bench
2}, pinned to revision
\texttt{2fd12b88aafdd04a52c298e3940bcb189f9766d6}. This is a real shell
and filesystem task using the original deterministic log generator and
unmodified official Python tests. Its Apache 2.0 license, author
attribution, source hashes, and evaluation-only canary are retained. It
is not one of the synthetic workflow worlds and never enters their
training exports.

The
\href{../../artifacts/public-benchmark/native-tools-v2-final/report.json}{final
native-tool receipt} records controller revision
\texttt{4b753e319fd013790524ae440d1abb9ecae9f92a} with
\texttt{source\_dirty:\ false}:

{\def\LTcaptype{none} % do not increment counter
\begin{longtable}[]{@{}
  >{\raggedright\arraybackslash}p{(\linewidth - 8\tabcolsep) * \real{0.1667}}
  >{\raggedright\arraybackslash}p{(\linewidth - 8\tabcolsep) * \real{0.1667}}
  >{\raggedleft\arraybackslash}p{(\linewidth - 8\tabcolsep) * \real{0.2222}}
  >{\raggedleft\arraybackslash}p{(\linewidth - 8\tabcolsep) * \real{0.2222}}
  >{\raggedleft\arraybackslash}p{(\linewidth - 8\tabcolsep) * \real{0.2222}}@{}}
\toprule\noalign{}
\begin{minipage}[b]{\linewidth}\raggedright
Attempt
\end{minipage} & \begin{minipage}[b]{\linewidth}\raggedright
Final official grade
\end{minipage} & \begin{minipage}[b]{\linewidth}\raggedleft
Commands executed
\end{minipage} & \begin{minipage}[b]{\linewidth}\raggedleft
Model calls
\end{minipage} & \begin{minipage}[b]{\linewidth}\raggedleft
Wall seconds
\end{minipage} \\
\midrule\noalign{}
\endhead
\bottomrule\noalign{}
\endlastfoot
Official reference algorithm & Pass, two tests & 1 & 0 & 3.42 \\
Deliberately wrong CSV & Fail, one of two tests & 1 & 0 & 0.67 \\
Native-tool local Qwen3 8B & Fail, empty CSV & 8 & 8 & 18.27 \\
\end{longtable}
}

The model was the already-installed Ollama \texttt{qwen3:8b}, Q4\_K\_M.
Its exact model digest is:

\begin{Shaded}
\begin{Highlighting}[]
\NormalTok{500a1f067a9f782620b40bee6f7b0c89e17ae61f686b92c24933e4ca4b2b8b41}
\end{Highlighting}
\end{Shaded}

No model download or paid call occurred. The final native-tool protocol
received eight valid tool calls and executed eight shell commands. The
attempt exhausted its eight-call budget with no execution exception but
left an empty CSV, which failed the official tests. This is an observed
bounded task failure, not a malformed-action failure.

Earlier JSON-v1 protocol evidence remains separate. Its
\href{../../artifacts/public-benchmark/attempts/corrected-harness-second-attempt/report.json}{corrected-harness
attempt} made eight local calls over 141.36 seconds; all responses ended
with \texttt{done\_reason:\ length} and an unterminated JSON command
string. No action reached the shell. That is an action-protocol and
output-budget failure. The earlier one-request harness receipt admits
its raw response was not captured. Neither missing provenance nor a
successful action is reconstructed. These are repeated attempts on the
same task under different controller protocols, not independent task
samples or evidence of changed model weights.

Agents run as a nonroot user with no container network, read-only root
filesystem, dropped capabilities, no privilege escalation,
CPU/memory/process limits, and a single writable CSV output mount. File
size is limited to 256 KiB, core dumps are disabled, and Docker
container logging is disabled. A real 512-KiB write probe failed at the
256-KiB kernel limit; the final receipt records agent and verifier
limits and no remaining episode containers. Agents receive neither
repository nor tests. The controller removes the agent, checks the
bounded regular output artifact, then grades in a fresh container with
tests mounted read-only. This is a concrete artifact handoff and
evaluator boundary, not a proof against host-kernel compromise.

The official solution's package-installation bootstrap is omitted, and
pytest 8.4.1 is installed during the verifier image build rather than
downloaded by the official test shell script during execution. Image
builds require network access; episodes and grading do not. The resolved
Python base digest and agent/verifier image IDs are recorded. These
disclosed adaptations are not the Terminal-Bench leaderboard protocol,
and a single task with one bounded model attempt cannot establish
general competence. \href{../public-benchmark.md}{Exact execution
protocol}.

\subsubsection{What remains unproved}\label{what-remains-unproved}

The final foundation verification ran 82 automated backend and
public-runner tests successfully, including envelope-forgery
regressions. The in-progress ATIF work is not part of this foundation
claim. One upstream Starlette test-client deprecation warning remains.
The browser receipt records 35 connected UI checks, no browser runtime
errors, desktop and 390-pixel mobile evidence, real matched
preference/grouped downloads, and cancellation through the UI. These
checks support the delivered paths, not exhaustive accessibility or
adversarial security coverage.
\href{../../artifacts/qa/ui-report.json}{Browser evidence}.

No measured result establishes production application parity,
human-calibrated task validity, task-family generalization, actual
SFT/DPO/GRPO learning, a training improvement, hosted-provider
performance, cloud isolation, multi-tenant security, or superiority to
Fleet, Proximal, Harbor, or another platform. The roadmap names the
required evidence rather than presenting those properties as delivered.

\end{document}
